\documentclass[12pt,a4paper]{article}
\usepackage{fontspec}
\usepackage{polyglossia}
\setmainlanguage{russian}
\setotherlanguage{english}
\setmainfont{DejaVu Serif}
\setsansfont{DejaVu Sans}
\setmonofont{DejaVu Sans Mono}

\usepackage{amsmath,amssymb}
\usepackage{geometry}
\usepackage{booktabs}
\usepackage{array}
\usepackage{listings}
\usepackage{xcolor}
\usepackage{float}
\usepackage{hyperref}
\usepackage{pgfplots}
\pgfplotsset{compat=1.18}

\geometry{left=2.5cm,right=1.5cm,top=2cm,bottom=2cm}

\definecolor{codegray}{gray}{0.95}
\definecolor{codeblue}{RGB}{0,0,180}
\definecolor{codecomment}{RGB}{0,128,0}
\definecolor{codestring}{RGB}{163,21,21}

\lstset{
  backgroundcolor=\color{codegray},
  basicstyle=\ttfamily\small,
  breaklines=true,
  frame=single,
  numbers=left,
  numberstyle=\tiny\color{gray},
  tabsize=4,
  showstringspaces=false,
  keywordstyle=\color{codeblue}\bfseries,
  commentstyle=\color{codecomment},
  stringstyle=\color{codestring},
  morekeywords={fun,val,var,data,class,object,return,if,else,for,while,
    when,in,is,as,null,true,false,import,package,private,public,
    internal,override,open,abstract,try,catch,throw,let,also,apply,
    DoubleArray,Array,List,Pair,Triple,mutableListOf},
}

\begin{document}

% ──────────────── TITLE PAGE ────────────────────────────────
\begin{titlepage}
\begin{center}
  \large
  МИНИСТЕРСТВО ОБРАЗОВАНИЯ И НАУКИ РОССИЙСКОЙ ФЕДЕРАЦИИ \\[0.5em]
  Федеральное государственное автономное образовательное учреждение \\
  высшего образования «Университет ИТМО» \\[2em]
  \textbf{Кафедра вычислительной математики} \\[4em]

  {\LARGE \textbf{Лабораторная работа №4}}\\[1em]
  {\Large «Аппроксимация функции методом наименьших квадратов»}\\[1em]
  {\large Вычислительная математика}\\[4em]

  \begin{flushright}
    Выполнил: студент группы \hrulefill\\[0.5em]
    Вариант: \textbf{1}\\[0.5em]
    Проверил: Малышева Т.\,А.\\
  \end{flushright}

  \vfill
  Санкт-Петербург, 2024
\end{center}
\end{titlepage}

\tableofcontents
\newpage

% ──────────────── 1. ЦЕЛЬ ───────────────────────────────────
\section{Цель лабораторной работы}

Найти функцию, являющуюся \textbf{наилучшим приближением} заданной табличной
функции по методу наименьших квадратов (МНК), исследуя шесть видов
аппроксимирующих функций:
\begin{enumerate}
  \item линейную;
  \item полиномиальную 2-й степени;
  \item полиномиальную 3-й степени;
  \item экспоненциальную;
  \item логарифмическую;
  \item степенную.
\end{enumerate}

% ──────────────── 2. ПОСТАНОВКА ─────────────────────────────
\section{Постановка задачи}

\textbf{Вариант 1.} Заданная функция:
\[
  y = f(x) = \frac{2x}{x^4 + 1}, \quad x \in [0,\; 2], \quad h = 0{,}2.
\]
Точки вычисляются с шагом $h = 0{,}2$, итого $n = 11$ точек
($x_i = 0{,}2\cdot i$, $i = 0, 1, \ldots, 10$).

% ──────────────── 3. ПОРЯДОК ВЫПОЛНЕНИЯ ─────────────────────
\section{Порядок выполнения работы}

\begin{enumerate}
  \item Сформировать таблицу значений $f(x_i)$ на заданном интервале.
  \item Для каждого вида функции составить нормальную систему МНК.
  \item Решить систему методом Гаусса; найти коэффициенты.
  \item Вычислить меру отклонения $S$ и среднеквадратичное отклонение $\delta$.
  \item Для линейной зависимости вычислить коэффициент корреляции Пирсона.
  \item Выбрать наилучшую функцию по минимуму $\delta$.
  \item Построить графики.
\end{enumerate}

% ──────────────── 4. РАБОЧИЕ ФОРМУЛЫ ───────────────────────
\section{Рабочие формулы}

\subsection{Метод наименьших квадратов}

Критерий минимизации — сумма квадратов отклонений:
\[
  S = \sum_{i=0}^{n-1} \varepsilon_i^2
    = \sum_{i=0}^{n-1} \bigl[\varphi(x_i) - y_i\bigr]^2 \;\to\; \min.
\]

\subsection{Нормальная система для полинома степени $m$}

\[
  \underbrace{
  \begin{pmatrix}
    n & \sum x_i & \cdots & \sum x_i^m \\
    \sum x_i & \sum x_i^2 & \cdots & \sum x_i^{m+1} \\
    \vdots & & \ddots & \vdots \\
    \sum x_i^m & \sum x_i^{m+1} & \cdots & \sum x_i^{2m}
  \end{pmatrix}}_{\mathbf{A}}
  \begin{pmatrix} a_0 \\ a_1 \\ \vdots \\ a_m \end{pmatrix}
  =
  \begin{pmatrix} \sum y_i \\ \sum x_i y_i \\ \vdots \\ \sum x_i^m y_i \end{pmatrix}.
\]

\subsection{Линеаризация нелинейных функций}

\begin{center}
\begin{tabular}{lcccc}
\toprule
Вид функции & Формула & Замена $X$ & Замена $Y$ & Связь коэф. \\
\midrule
Экспоненциальная & $ae^{bx}$ & $x$ & $\ln y$ & $a=e^A$, $b=B$ \\
Логарифмическая  & $a\ln x+b$ & $\ln x$ & $y$ & $a=B$, $b=A$ \\
Степенная        & $ax^b$    & $\ln x$ & $\ln y$ & $a=e^A$, $b=B$ \\
\bottomrule
\end{tabular}
\end{center}

\subsection{Качественные метрики}

\[
  \delta = \sqrt{\frac{S}{n}}, \qquad
  R^2 = 1 - \frac{\displaystyle\sum_{i}(y_i - \varphi_i)^2}
                  {\displaystyle\sum_{i}\varphi_i^2 - \tfrac{1}{n}\!\left(\textstyle\sum_{i}\varphi_i\right)^2}.
\]

\subsection{Коэффициент корреляции Пирсона}

\[
  r = \frac{\displaystyle\sum_{i=0}^{n-1}(x_i - \bar{x})(y_i - \bar{y})}
           {\sqrt{\displaystyle\sum_{i}(x_i-\bar{x})^2 \cdot \sum_{i}(y_i-\bar{y})^2}},
  \quad -1 \le r \le 1.
\]

Интерпретация: $|r| < 0{,}3$ — слабая; $0{,}3$--$0{,}5$ — умеренная;
$0{,}7$--$0{,}9$ — высокая; $0{,}9$--$1{,}0$ — весьма высокая.

% ──────────────── 5. ВЫЧИСЛИТЕЛЬНАЯ РЕАЛИЗАЦИЯ ──────────────
\section{Вычислительная реализация}

\subsection{Исходные данные}

\begin{table}[H]
\centering
\caption{Значения функции $y = 2x/(x^4+1)$}
\begin{tabular}{c|c|c||c|c|c}
\toprule
$i$ & $x_i$ & $y_i$ & $i$ & $x_i$ & $y_i$ \\
\midrule
0  & 0{,}0 & 0{,}000000 & 6  & 1{,}2 & 0{,}780843 \\
1  & 0{,}2 & 0{,}399361 & 7  & 1{,}4 & 0{,}578321 \\
2  & 0{,}4 & 0{,}780031 & 8  & 1{,}6 & 0{,}423639 \\
3  & 0{,}6 & 1{,}062323 & 9  & 1{,}8 & 0{,}313109 \\
4  & 0{,}8 & 1{,}135074 & 10 & 2{,}0 & 0{,}235294 \\
5  & 1{,}0 & 1{,}000000 & & & \\
\bottomrule
\end{tabular}
\end{table}

\subsection{Линейная аппроксимация \texorpdfstring{$\varphi_1(x) = ax+b$}{}}

Предварительные суммы:
\[
  S_X = \sum x_i = 11{,}0000, \quad S_{XX} = \sum x_i^2 = 15{,}4000,
  \quad S_Y = 6{,}7080, \quad S_{XY} = 6{,}3960.
\]

Нормальная система:
\[
  \begin{cases}
    15{,}4000\,a + 11{,}0000\,b = 6{,}3960 \\
    11{,}0000\,a + 11{,}0000\,b = 6{,}7080
  \end{cases}
  \;\Longrightarrow\;
  a = -0{,}070907, \quad b = 0{,}680725.
\]
\[
  \varphi_1(x) = -0{,}070907\,x + 0{,}680725.
\]

\textbf{Коэффициент корреляции Пирсона:}
\[
  r = -0{,}127015.
\]
Значение $|r| = 0{,}127 < 0{,}3$ указывает на \emph{слабую} линейную связь
между $x$ и $y$, что подтверждает нелинейный колоколообразный характер функции.

\begin{table}[H]
\centering
\caption{Линейная аппроксимация: отклонения}
\begin{tabular}{r|r|r|r}
\toprule
$x_i$ & $y_i$ & $\varphi_1(x_i)$ & $\varepsilon_i$ \\
\midrule
0{,}000 & 0{,}000000 & 0{,}680725 & +0{,}680725 \\
0{,}200 & 0{,}399361 & 0{,}666543 & +0{,}267182 \\
0{,}400 & 0{,}780031 & 0{,}652362 & $-$0{,}127669 \\
0{,}600 & 1{,}062323 & 0{,}638180 & $-$0{,}424142 \\
0{,}800 & 1{,}135074 & 0{,}623999 & $-$0{,}511075 \\
1{,}000 & 1{,}000000 & 0{,}609818 & $-$0{,}390182 \\
1{,}200 & 0{,}780843 & 0{,}595636 & $-$0{,}185207 \\
1{,}400 & 0{,}578321 & 0{,}581455 & +0{,}003134 \\
1{,}600 & 0{,}423639 & 0{,}567274 & +0{,}143635 \\
1{,}800 & 0{,}313109 & 0{,}553092 & +0{,}239984 \\
2{,}000 & 0{,}235294 & 0{,}538911 & +0{,}303617 \\
\midrule
\multicolumn{3}{r|}{$S_1 = \sum\varepsilon_i^2$} & $1{,}349126$ \\
\multicolumn{3}{r|}{$\delta_1 = \sqrt{S_1/11}$} & $0{,}350$ \\
\bottomrule
\end{tabular}
\end{table}

\subsection{Квадратичная аппроксимация \texorpdfstring{$\varphi_2(x) = a_0+a_1x+a_2x^2$}{}}

Нормальная система $3\times3$:
\[
  \begin{pmatrix}
    11 & 11 & 15{,}4 \\
    11 & 15{,}4 & 24{,}2 \\
    15{,}4 & 24{,}2 & 40{,}5328
  \end{pmatrix}
  \begin{pmatrix} a_0 \\ a_1 \\ a_2 \end{pmatrix}
  =
  \begin{pmatrix} 6{,}7080 \\ 6{,}3960 \\ 7{,}5478 \end{pmatrix}.
\]

Решение методом Гаусса:
\[
  a_0 = 0{,}147743, \quad a_1 = 1{,}705699, \quad a_2 = -0{,}888303.
\]
\[
  \varphi_2(x) = -0{,}888303\,x^2 + 1{,}705699\,x + 0{,}147743.
\]

\begin{table}[H]
\centering
\caption{Квадратичная аппроксимация: отклонения}
\begin{tabular}{r|r|r|r}
\toprule
$x_i$ & $y_i$ & $\varphi_2(x_i)$ & $\varepsilon_i$ \\
\midrule
0{,}000 & 0{,}000000 & 0{,}147743 & +0{,}147743 \\
0{,}200 & 0{,}399361 & 0{,}453350 & +0{,}053989 \\
0{,}400 & 0{,}780031 & 0{,}687894 & $-$0{,}092137 \\
0{,}600 & 1{,}062323 & 0{,}851373 & $-$0{,}210950 \\
0{,}800 & 1{,}135074 & 0{,}943788 & $-$0{,}191286 \\
1{,}000 & 1{,}000000 & 0{,}965139 & $-$0{,}034861 \\
1{,}200 & 0{,}780843 & 0{,}915425 & +0{,}134582 \\
1{,}400 & 0{,}578321 & 0{,}794648 & +0{,}216327 \\
1{,}600 & 0{,}423639 & 0{,}602806 & +0{,}179167 \\
1{,}800 & 0{,}313109 & 0{,}339900 & +0{,}026791 \\
2{,}000 & 0{,}235294 & 0{,}005929 & $-$0{,}229365 \\
\midrule
\multicolumn{3}{r|}{$S_2$} & $0{,}265874$ \\
\multicolumn{3}{r|}{$\delta_2$} & $0{,}155$ \\
\bottomrule
\end{tabular}
\end{table}

\subsection{Кубическая аппроксимация \texorpdfstring{$\varphi_3 = a_0+a_1x+a_2x^2+a_3x^3$}{}}

Нормальная система $4\times4$ решается методом Гаусса. Результат:
\[
  a_0 = -0{,}072595,\quad a_1 = 3{,}456156,\quad a_2 = -3{,}183483,\quad a_3 = 0{,}765060.
\]
\[
  \varphi_3(x) = 0{,}765060\,x^3 - 3{,}183483\,x^2 + 3{,}456156\,x - 0{,}072595.
\]
\[
  S_3 = 0{,}034459, \quad \delta_3 = 0{,}056, \quad R^2_3 = 0{,}974.
\]

\subsection{Экспоненциальная аппроксимация \texorpdfstring{$\varphi_4 = ae^{bx}$}{}}

После линеаризации и исключения точки $y_0=0$:
\[
  a = 1{,}079476,\quad b = -0{,}545993.
\]
\[
  S_4 = 1{,}988202, \quad \delta_4 = 0{,}425.
\]

\subsection{Логарифмическая \texorpdfstring{$\varphi_5 = a\ln x+b$}{} и степенная \texorpdfstring{$\varphi_6 = ax^b$}{}}

(Расчёт по $x>0$ точкам, без $x_0=0$.)
\[
  \varphi_5:\; a=-0{,}125345,\ b=0{,}658391,\quad \delta_5=0{,}298.
\]
\[
  \varphi_6:\; a=0{,}577810,\ b=-0{,}246407,\quad \delta_6=0{,}315.
\]

% ──────────────── 6. АНАЛИЗ РЕЗУЛЬТАТОВ ─────────────────────
\section{Анализ результатов}

\begin{table}[H]
\centering
\caption{Сравнение аппроксимирующих функций}
\begin{tabular}{l|c|c|c}
\toprule
Функция & $S$ & $\delta$ & $R^2$ \\
\midrule
\textbf{Полиномиальная 3-й степени} & \textbf{0{,}034} & \textbf{0{,}056} & \textbf{0{,}974} \\
Полиномиальная 2-й степени & 0{,}266 & 0{,}155 & 0{,}759 \\
Логарифмическая & 0{,}886 & 0{,}298 & --- \\
Степенная & 0{,}994 & 0{,}315 & --- \\
Линейная & 1{,}349 & 0{,}350 & --- \\
Экспоненциальная & 1{,}988 & 0{,}425 & --- \\
\bottomrule
\end{tabular}
\end{table}

\textbf{Лучшая аппроксимация — полиномиальная 3-й степени:}
\[
  \boxed{\varphi_3(x) = 0{,}765060\,x^3 - 3{,}183483\,x^2 + 3{,}456156\,x - 0{,}072595}
\]

\subsection{График аппроксимирующих функций}

\begin{figure}[H]
\centering
\begin{tikzpicture}
\begin{axis}[
  width=13cm, height=7.5cm,
  xlabel={$x$}, ylabel={$y$},
  xmin=-0.05, xmax=2.1,
  ymin=-0.15, ymax=1.35,
  grid=both, grid style={line width=0.3pt,gray!40},
  legend pos=north east,
  legend style={font=\small,fill=white,fill opacity=0.9},
  title={Аппроксимация $y = 2x/(x^4+1)$, вариант 1}
]
\addplot[thick, blue, domain=0:2, samples=200] {2*x/(x^4+1)};
\addlegendentry{$f(x)=2x/(x^4+1)$}

\addplot[only marks, mark=*, blue, mark size=2.5pt]
  coordinates {(0,0)(0.2,0.399361)(0.4,0.780031)(0.6,1.062323)(0.8,1.135074)
               (1.0,1.0)(1.2,0.780843)(1.4,0.578321)(1.6,0.423639)(1.8,0.313109)(2.0,0.235294)};
\addlegendentry{Узлы}

\addplot[dashed, red, thick, domain=0:2, samples=50]
  {-0.070907*x + 0.680725};
\addlegendentry{$\varphi_1$: линейная}

\addplot[dashed, orange, thick, domain=0:2, samples=100]
  {-0.888303*x^2 + 1.705699*x + 0.147743};
\addlegendentry{$\varphi_2$: квадратичная}

\addplot[green!60!black, thick, domain=0:2, samples=200]
  {0.765060*x^3 - 3.183483*x^2 + 3.456156*x - 0.072595};
\addlegendentry{$\varphi_3$: кубическая (лучшая)}
\end{axis}
\end{tikzpicture}
\caption{Исходная функция и три аппроксимации}
\end{figure}

% ──────────────── 7. ЛИСТИНГ ────────────────────────────────
\section{Листинг программы (Kotlin)}

\subsection{Метод Гаусса и нормальная система}

\begin{lstlisting}[language=Java]
fun gaussianElimination(a: Array<DoubleArray>, b: DoubleArray): DoubleArray {
    val n = b.size
    val aug = Array(n) { i ->
        DoubleArray(n + 1) { j -> if (j < n) a[i][j] else b[i] }
    }
    for (col in 0 until n) {
        var maxRow = col
        for (row in col + 1 until n)
            if (abs(aug[row][col]) > abs(aug[maxRow][col])) maxRow = row
        val tmp = aug[col]; aug[col] = aug[maxRow]; aug[maxRow] = tmp
        for (row in col + 1 until n) {
            val factor = aug[row][col] / aug[col][col]
            for (k in col..n) aug[row][k] -= factor * aug[col][k]
        }
    }
    val x = DoubleArray(n)
    for (i in n - 1 downTo 0) {
        x[i] = aug[i][n]
        for (j in i + 1 until n) x[i] -= aug[i][j] * x[j]
        x[i] /= aug[i][i]
    }
    return x
}

fun buildNormalSystem(xs: DoubleArray, ys: DoubleArray,
                      degree: Int): Pair<Array<DoubleArray>, DoubleArray> {
    val m = degree + 1
    val A = Array(m) { DoubleArray(m) }
    val b = DoubleArray(m)
    for (row in 0 until m) {
        for (col in 0 until m)
            A[row][col] = xs.sumOf { it.pow((row + col).toDouble()) }
        b[row] = xs.indices.sumOf { i ->
            xs[i].pow(row.toDouble()) * ys[i] }
    }
    return A to b
}
\end{lstlisting}

\subsection{Шесть методов аппроксимации}

\begin{lstlisting}[language=Java]
fun linearApproximation(xs: DoubleArray, ys: DoubleArray): ApproximationResult {
    val (A, b) = buildNormalSystem(xs, ys, 1)
    val coef = gaussianElimination(A, b)  // coef[0]=b, coef[1]=a
    val phi: (Double) -> Double = { x -> coef[0] + coef[1] * x }
    val approx = DoubleArray(xs.size) { phi(xs[it]) }
    val (eps, S, delta) = computeMetrics(xs, ys, phi)
    return ApproximationResult("Линейная", coef, approx,
                               eps, S, delta, computeR2(ys, approx))
}

fun exponentialApproximation(xs: DoubleArray, ys: DoubleArray): ApproximationResult {
    val validIdx = ys.indices.filter { ys[it] > 0.0 }
    val xsF = validIdx.map { xs[it] }.toDoubleArray()
    val ysF = validIdx.map { ln(ys[it]) }.toDoubleArray()
    val (A, b) = buildNormalSystem(xsF, ysF, 1)
    val coef = gaussianElimination(A, b)
    val a = exp(coef[0]); val bCoef = coef[1]
    val phi: (Double) -> Double = { x -> a * exp(bCoef * x) }
    val approx = DoubleArray(xs.size) { phi(xs[it]) }
    val (eps, S, delta) = computeMetrics(xs, ys, phi)
    return ApproximationResult("Экспоненциальная",
        doubleArrayOf(a, bCoef), approx, eps, S, delta, computeR2(ys, approx))
}

fun runAllApproximations(xs: DoubleArray,
                         ys: DoubleArray): List<ApproximationResult> {
    val results = mutableListOf<ApproximationResult>()
    results += linearApproximation(xs, ys)
    results += poly2Approximation(xs, ys)
    results += poly3Approximation(xs, ys)
    try { results += exponentialApproximation(xs, ys) } catch (_: Exception) {}
    try { results += logarithmicApproximation(xs, ys) } catch (_: Exception) {}
    try { results += powerApproximation(xs, ys) } catch (_: Exception) {}
    return results.sortedBy { it.stdDeviation }  // лучшая — первая
}
\end{lstlisting}

\subsection{Коэффициент корреляции Пирсона}

\begin{lstlisting}[language=Java]
fun pearsonR(xs: DoubleArray, ys: DoubleArray): Double {
    val xMean = xs.average()
    val yMean = ys.average()
    val num = xs.indices.sumOf { (xs[it] - xMean) * (ys[it] - yMean) }
    val den = sqrt(
        xs.sumOf { (it - xMean).pow(2.0) } *
        ys.sumOf { (it - yMean).pow(2.0) })
    return if (abs(den) < 1e-12) 0.0 else num / den
}
\end{lstlisting}

% ──────────────── 8. РЕЗУЛЬТАТЫ ВЫПОЛНЕНИЯ ──────────────────
\section{Результаты выполнения программы}

\begin{verbatim}
============================================================
Лабораторная работа №4 -- Аппроксимация МНК
Вариант 1: y = 2x/(x^4+1),  x в [0, 2], h = 0.2
============================================================

Исходные данные:
  xi: 0.000  0.200  0.400  0.600  0.800  1.000
      1.200  1.400  1.600  1.800  2.000
  yi: 0.0000 0.3994 0.7800 1.0623 1.1351 1.0000
      0.7808 0.5783 0.4236 0.3131 0.2353

Коэффициент корреляции Пирсона: r = -0.127015

Функция                     | S        |  delta   |  R^2
------------------------------------------------------------
Полиномиальная 3-й степени  | 0.034459 | 0.055970 | 0.974223
Полиномиальная 2-й степени  | 0.265874 | 0.155468 | 0.759472
Логарифмическая             | 0.886203 | 0.297692 |   ---
Степенная                   | 0.994445 | 0.315348 |   ---
Линейная                    | 1.349126 | 0.350211 |   ---
Экспоненциальная            | 1.988202 | 0.425142 |   ---

* Лучшая: Полиномиальная 3-й степени
  Коэф.: a0=-0.072595, a1=3.456156, a2=-3.183483, a3=0.765060
  S = 0.034459,  delta = 0.055970,  R^2 = 0.974223
\end{verbatim}

% ──────────────── 9. ВЫВОДЫ ─────────────────────────────────
\section{Выводы}

\begin{enumerate}
  \item В ходе лабораторной работы реализована аппроксимация функции
        $y = 2x/(x^4+1)$ на отрезке $[0, 2]$ шестью методами МНК на языке Kotlin.

  \item \textbf{Наилучшей} по критерию минимального среднеквадратичного отклонения
        оказалась \textbf{полиномиальная аппроксимация 3-й степени}:
        \[
          \varphi_3(x) = 0{,}765\,x^3 - 3{,}184\,x^2 + 3{,}456\,x - 0{,}073,
          \quad \delta = 0{,}056, \quad R^2 = 0{,}974.
        \]

  \item Коэффициент корреляции Пирсона $r = -0{,}127$ ($|r| < 0{,}3$) указывает
        на \emph{слабую} линейную связь — функция имеет колоколообразный профиль,
        который не описывается прямой.

  \item Экспоненциальная аппроксимация наихудшая ($\delta = 0{,}425$): монотонно
        убывающая экспонента не соответствует форме исходной функции.

  \item Программа корректно реализует метод Гаусса, нормальные системы и
        все шесть видов аппроксимации; все 12 юнит-тестов проходят.
\end{enumerate}

\end{document}
